\section{النموذج الديناميكي العام وفق صياغة لاغرانج}

سيتم اشتقاق النموذج الديناميكي للذراع الروبوتية باستخدام طريقة لاغرانج 
\textenglish{(Lagrangian Method)}، وهي طريقة تعتمد على توصيف النظام من خلال طاقته الحركية وطاقته الكامنة، ثم اشتقاق معادلات الحركة الناتجة عن ذلك. تُعد هذه الطريقة مناسبة للأنظمة الروبوتية متعددة المفاصل، لأنها تسمح بالحصول على نموذج ديناميكي عام يربط بين العزوم المطبقة على المفاصل والحركة الناتجة عنها.

يعطى الشكل العام لمعادلات الحركة الناتجة عن تطبيق طريقة لاغرانج على ذراع روبوتية مكونة من $n$ درجة حرية كما يلي:

\begin{equation}
T_i =
\sum_{j=1}^{n} D_{ij}\ddot{q}_j
+
I_{i(\mathrm{act})}\ddot{q}_i
+
\sum_{j=1}^{n}\sum_{k=1}^{n}
D_{ijk}\dot{q}_j\dot{q}_k
+
D_i
\end{equation}

حيث إن:

\begin{itemize}
    \item $T_i$ يمثل العزم المطبق على المفصل $i$.
    \item $q_i$ يمثل متحول المفصل $i$.
    \item $\dot{q}_i$ يمثل سرعة المفصل $i$.
    \item $\ddot{q}_i$ يمثل تسارع المفصل $i$.
    \item $I_{i(\mathrm{act})}$ يمثل عزم القصور الذاتي للمشغل المرتبط بالمفصل $i$.
\end{itemize}

تتكون المعادلة السابقة من ثلاثة حدود رئيسية: حد العطالة، وحد قوى كوريوليس والقوى الطاردة المركزية، وحد الجاذبية. وتحسب هذه الحدود كما يلي:

\subsection{مصفوفة العطالة}

تمثل $D_{ij}$ عناصر مصفوفة العطالة، وهي المسؤولة عن تأثير كتل الوصلات وعزوم قصورها الذاتي على حركة المفاصل، وتعطى بالعلاقة:

\begin{equation}
D_{ij} =
\sum_{p=\max(i,j)}^{n}
\operatorname{Trace}
\left(
U_{pj} J_p U_{pi}^{T}
\right)
\end{equation}

حيث تمثل $J_p$ مصفوفة القصور الذاتي للوصلة $p$، بينما تمثل $U_{pj}$ مشتق مصفوفة التحويل الخاصة بالوصلة $p$ بالنسبة لمتحول المفصل $q_j$.

\subsection{مصفوفة قوى كوريوليس والقوى الطاردة المركزية}

تمثل $D_{ijk}$ معاملات قوى كوريوليس والقوى الطاردة المركزية، والتي تظهر نتيجة حركة المفاصل بسرعات زاوية غير معدومة، وتعطى بالعلاقة:

\begin{equation}
D_{ijk} =
\sum_{p=\max(i,j,k)}^{n}
\operatorname{Trace}
\left(
U_{pjk} J_p U_{pi}^{T}
\right)
\end{equation}

حيث تمثل $U_{pjk}$ المشتقات الجزئية من المرتبة الثانية لمصفوفات التحويل بالنسبة لمتحولات المفاصل.

\subsection{مصفوفة الجاذبية}

يمثل $D_i$ تأثير الجاذبية على المفصل $i$، ويحسب بالعلاقة:

\begin{equation}
D_i =
\sum_{p=i}^{n}
- m_p g^{T} U_{pi} \bar{r}_p
\end{equation}

حيث إن $m_p$ تمثل كتلة الوصلة $p$، و $g$ يمثل شعاع تسارع الجاذبية، أما $\bar{r}_p$ فيمثل شعاع مركز كتلة الوصلة بالنسبة إلى إطارها المحلي.

وبذلك نحصل على الشكل العام للنموذج الديناميكي للذراع الروبوتية باستخدام طريقة لاغرانج، والذي يصف العلاقة بين عزوم المفاصل من جهة، وتسارعات وسرعات المفاصل وتأثير الجاذبية من جهة أخرى.


